\scanpage{79}
% archive foliation: 81
plongés dans $X''$ normalisé de $X'''$ \deleted{dont}\inserted{duquel} on
part, et \deleted{on} trouve des antiiso entre
leurs bords \deleted{de ces} de $D_i',D_i''$ ---
d'où possibilité de les recoller pour
trouver un $X$ surface holo (à bord,
si on a aussi enlevé des disques ouv.
$\mathring D_\alpha^{\unsure{*}}$ correspondant aux pts marqués $x_\alpha$,
sinon une surface à pts marqués),
découpée par les cercles $\Gamma_i$ qui ne
sont \deleted{pas}\inserted{peut-être pas} \unsure{néc.} géodésiques
($??$ \deleted{\unsure{on a}}\inserted{à}
examiner de plus près) mais définissant
des cercles géodésiques uniques,
par les principes de Thurston.

Il faudrait voir que le résultat
obtenu ne dépend pas du choix des
paires $(D_i',D_i'')$ mut.\textsuperscript{t} polaires, il y
a un argument heuristique pour
cela tout à fait vaseux, consistant
à dire qu'une petite variation de \unsure{celle|cette},
